% problem-000088 4 聚合物合成与热力学
\section*{4. 聚合物合成与热力学}
（图：）
A
B

\subsection*{4-1}
⽐较A与B开环聚合时热⼒学驱动⼒的⼤⼩ $\setminus _ { \circ }$ 。简述原因。

\subsection*{4-2}
在有机弱碱的催化作⽤下，苄硫醇作为引发剂，B在室温下可发⽣活性开环聚合，形成聚合产物。画出产物的结构式。

\subsection*{4-3}
聚合 $( \mathbf { B } _ { n - 1 } ^ { \prime }   \mathbf { B } _ { n } ^ { \prime } )$ 过程中， $\Delta H ^ { \ominus } = - 1 5 . 6$ kJ $\mathrm { m o l ^ { - 1 } }$ $\Delta S ^ { \ominus } = - 4 0 . 4 \mathrm { J } \mathrm { m o l } ^ { - 1 } \mathrm { K } _ { \circ } ^ { - 1 }$ 计算室温(298 K)下反应的平衡常数 $K ^ { \ominus }$ ，设反应达平衡时单体浓度 $[ \mathbf { B } ] _ { \mathrm { e q } } / c ^ { \epsilon }$ ⊖等于 $1 / K$ ，若起始单体浓度为2.00 mol L−1，引发剂浓度为0.0100 mol L−1，计算达平衡时产物的平均聚合度 $\dot { \boldsymbol { n } } _ { < }$ 。

\subsection*{4-4}
提⾼单体B开环聚合转化率的⽅法有：

a) 升⾼温度； b) 降低温度； c) 增加B的起始浓度； d) 延⻓反应时间
